\documentclass[10pt]{article}
% Math Packages
\usepackage{amsmath, mathtools}
\usepackage{amssymb}
\usepackage{amsthm}
\usepackage{amsfonts}
\usepackage{bbm}
\usepackage{breqn}
\usepackage[margin=1in]{geometry}
\usepackage{graphicx}
\usepackage{tikz}
\usetikzlibrary{arrows.meta}
\usetikzlibrary{calc}
\usepackage{forest}
\usepackage{tikz-qtree}
\graphicspath{ {./images/} }
\usepackage{hyperref}
\usepackage[capitalize]{cleveref}
\usepackage[shortlabels]{enumitem}
\usetikzlibrary{arrows,matrix,positioning}
\usepackage{multicol}

% for the pipe symbol
\usepackage[T1]{fontenc}

% Citing theorems by name. (source: https://tex.stackexchange.com/questions/109843/cleveref-and-named-theorems)
\makeatletter
\newcommand{\ncref}[1]{\cref{#1}\mynameref{#1}{\csname r@#1\endcsname}}

\def\mynameref#1#2{%
  \begingroup
  \edef\@mytxt{#2}%
  \edef\@mytst{\expandafter\@thirdoffive\@mytxt}%
  \ifx\@mytst\empty\else
  \space(\nameref{#1})\fi
  \endgroup
}
\makeatother

% Colorful Notes
\usepackage{color} \definecolor{Red}{rgb}{1,0,0} \definecolor{Blue}{rgb}{0,0,1}
\definecolor{Purple}{rgb}{.5,0,.5} \def\red{\color{Red}} \def\blue{\color{Blue}}
\def\gray{\color{gray}} \def\purple{\color{Purple}}
\newcommand{\rnote}[1]{{\red [#1]}} % \rnote{foo} gives '[foo]' in red
\newcommand{\pnote}[1]{{\purple [#1]}} % \pnote{foo} gives '[foo]' in purple
\newcommand{\bnote}[1]{{\blue #1}} % \bnote{foo} gives 'foo' in blue
\newcommand{\gnote}[1]{{\gray #1}} % \gnote{foo} gives 'foo' in gray
\newcommand{\Max}[1]{{\purple [#1]}} % \bnote{foo} then 'foo' is blue


% Claim numbering (the counter restarts after each proof environment)
\newcounter{claimcount}
\setcounter{claimcount}{0}
\newenvironment{claim}{\refstepcounter{claimcount}\par\addvspace{\medskipamount}\noindent\textbf{Claim \arabic{claimcount}:}}{}
\usepackage{etoolbox}
\AtBeginEnvironment{proof}{\setcounter{claimcount}{0}}
\newenvironment{claimproof}{\par\addvspace{\medskipamount}\noindent\textit{Proof of Claim  \arabic{claimcount}.}}{\hfill\ensuremath{\qedsymbol} \tiny{Claim}

  \medskip}
% Add claim support to cleverref
\crefname{claimcount}{Claim}{Claims}


% Math Environments
\newtheorem{theorem}{Theorem}
\newtheorem{assumption}[theorem]{Assumption}
\newtheorem{lemma}[theorem]{Lemma}
\newtheorem{proposition}[theorem]{Proposition}
\newtheorem{corollary}[theorem]{Corollary}
\newtheorem{question}[theorem]{Question}
\theoremstyle{definition}
\newtheorem{definition}[theorem]{Definition}
\newtheorem{remark}[theorem]{Remark}
\newtheorem{example}[theorem]{Example}
\newtheorem{notation}[theorem]{Notation}
\newtheorem{problem}[theorem]{Problem}

% Matrices and Column Vectors. 
\usepackage{stackengine}
\setstackgap{L}{1.0\normalbaselineskip}
\usepackage{tabstackengine}
\setstackEOL{;}% row separator
\setstackTAB{,}% column separator
\setstacktabbedgap{1ex}% inter-column gap
\setstackgap{L}{1.5\normalbaselineskip}% inter-row baselineskip
\let\nmatrix\bracketMatrixstack  %Usage: \nmatrix{1,2,3\4,5,6}
\newcommand\cv[1]{\setstackEOL{,}\bracketMatrixstack{#1}} %usage: \cv{1,2,3}

% Custom Math Coqmmands
\newcommand{\vt}{\vskip 5mm} % vertical space
\newcommand{\fl}{\noindent\textbf} % first line
\newcommand{\Fl}{\vt\noindent\textbf} % first line with space above
\newcommand{\norm}[1]{\left\lVert#1\right\rVert} % norm
\newcommand{\pnorm}[1]{\left\lVert#1\right\rVert_p} % p-norm
\newcommand{\qnorm}[1]{\left\lVert#1\right\rVert_q} % q-norm
\newcommand{\1}[1]{\textbf{1}_{\left[#1\right]}} % indicator function 
\def\limn{\lim_{n\to\infty}} % shortcut for lim as n-> infinity
\def\sumn{\sum_{n=1}^{\infty}} % shortcut for sum from n=1 to infinity
\def\sumkn{\sum_{k=1}^{n}} % shortcut for sum from k=1 to n
\def\sumin{\sum_{i=1}^{n}} % shortcut for sum from i=1 to n
\def\SAs{\sigma\text{-algebras}} % shortcut for $\sigma$-algebras
\def\SA{\sigma\text{-algebra}} % shortcut for $\sigma$-algebra
\def\Ft{\mathcal{F}_t} % time-indexed sigma-algebra (t)
\def\Fs{\mathcal{F}_s} % time-indexed sigma-algebra (s)
\def\F{\mathcal{F}} % sigma-algebra
\def\G{\mathcal{G}} % sigma-algebra
\def\R{\mathbb{R}} % Real numbers
\def\N{\mathbb{N}} % Natural numbers
\def\Z{\mathbb{Z}} % Integers
\def\E{\mathbb{E}} % Expectation
\def\P{\mathbb{P}} % Probability
\def\Q{\mathbb{Q}} % Q probability
\def\dist{\text{dist}} %Text 'dist' for things like 'dist(x,y)'
\newcommand{\indep}{\perp \!\!\! \perp}  %independence symbol
\def\Var{\mathrm{Var}} % Variance
\def\tr{\mathrm{tr}} % trace

% Brackets and Parentheses
\def\[{\left [}
    \def\]{\right ]}
% \def\({\left (}
%   \def\){\right )}



\usepackage{color}
\definecolor{Red}{rgb}{1,0,0}
\definecolor{Blue}{rgb}{0,0,1}
\definecolor{Purple}{rgb}{.5,0,.5}
\def\red{\color{Red}}
\def\blue{\color{Blue}}
\def\gray{\color{gray}}
\def\purple{\color{Purple}}
\definecolor{RoyalBlue}{cmyk}{1, 0.50, 0, 0}
\newcommand{\dempfcolor}[1]{{\color{RoyalBlue}#1}} 
\newcommand{\demph}[1]{\dempfcolor{{\sl #1}}}

% comment exactly one of the following line to show / hide the solutions
% \newcommand{\solution}[1]{{\purple #1}} % uncomment to show the solutions
\newcommand{\solution}[1]{} % uncomment to hide the solutions



\title{Lecture Notes for Math 372: \\Elementary Probability and Statistics}
\date{Last updated: \today}
% \author{mh}

\begin{document}
\begin{center}
  \section*{Math 372: Homework 04}
  \textit{Due Wednesday, Feb 11}
\end{center}

\begin{problem}[Exercise 2.3.34]

  \textbf{This problem was accidentally repeated from HW 3. You can ignore
    this problem.}
  
\end{problem}

\begin{problem}[Exercise 2.3.39]
  A box in a supply room contains 15 compact fluorescent lightbulbs, of which
  5 are rated 13-watt, 6 are rated 18-watt, and 4 are rated 23-watt. Suppose
  that three of these bulbs are randomly selected.
  \begin{enumerate}[(a)]
    \item What is the probability that exactly two of the selected bulbs are
    rated 23-watt?
    \item What is the probability that all three of the bulbs have the same
    rating?
    \item What is the probability that one bulb of each type is selected?
    \item If bulbs are selected one by one until a 23-watt bulb is obtained,
    what is the probability that it is necessary to examine at least 6 bulbs?
  \end{enumerate}
\end{problem}

\begin{problem}[Exercise 2.4.47]
  Consider randomly selecting a student at a large university. Let $A$ be the
  event that the selected student has a Visa card, let $B$ be the event that
  the student has a MasterCard, and let $C$ be the event that they have an
  American Express card. Suppose that
  \begin{gather}
    \P\left[A \right] = .6, \quad \P\left[B \right] = .4, \quad \P\left[C \right] = .2\\
    \P\left[A\cap B \right] = .3, \quad \P\left[A\cap C \right] = .15, \quad \P\left[B\cap C \right] = .1, \\
    \quad \text{and} \quad \P\left[A\cap B\cap C \right]=.08.
  \end{gather}
  \begin{enumerate}[(a)]
    \item What is the probability that the selected student has at least one
    of the three types of cards?
    \item What is the probability that the selected student has both a Visa
    card and a MasterCard, but not an American Express card?
    \item Calculate and interpret $\P\left[B\mid A \right]$ and $\P\left[A\mid
      B\right]$.
    \item If we learn that the selected student has an American Express card,
    what is the probability that she or he also has both a Visa card and a
    MasterCard?
    \item Given that the selected student has an American Express card, what
    is the probability that she or he has at leaset one of the other two types
    of cards?
  \end{enumerate}
\end{problem}

\begin{problem}[2.4.63]
  For customers purchasing a refrigerator at a certain appliance store, let
  $A$ be the event that the refrigerator was manufactured in the U.S., $B$ be
  the event that the refrigerator had an icemaker, and $C$ be the event that the
  customer purchased an extended warranty.

  Relevant probabilities are
  \begin{gather*}
    \P\left[A \right] = .75 \quad \P\left[B\mid A \right] = .9 \quad \P\left[B\mid A^{c} \right]=.8\\
    \P\left[C\mid A\cap B \right] = .8 \quad \P\left[C\mid A\cap B^{c} \right] = .6\\
    \P\left[C\mid A^{c}\cap B \right]=.7 \quad \P\left[C\mid A^{c}\cap B^{c} \right]=.3.
  \end{gather*}
  \begin{enumerate}[(a)]
    \item Construct a tree diagram consisting of first-, second-, and
    third-generation branches, and place an event label and
    appropriate probability next to each branch.
    \item Compute $\P\left[A\cap B\cap C \right]$.
    \item Compute $\P\left[B\cap C \right]$.
    \item Compute $\P\left[C \right]$.
    \item Compute $\P\left[A\mid B\cap C \right]$, the probability of a
    U.S.~purchase given that an icemaker and extended warranty are also
    purchased.
  \end{enumerate}
\end{problem}



\begin{problem}[Exercise 102 on p.93]
  A sample of 2026 students were asked about their eye color, with the
  following results:

  \begin{table}[h]
    \centering
    \begin{tabular}{|l|lllll|}
      \hline
       & \multicolumn{1}{l|}{\textbf{Blue}} & \multicolumn{1}{l|}{\textbf{Brown}} & \multicolumn{1}{l|}{\textbf{Green}} & \multicolumn{1}{l|}{\textbf{Hazel}} & \textbf{Total} \\ \hline
      \textbf{Male}   & \multicolumn{1}{l|}{370}           & \multicolumn{1}{l|}{352}            & \multicolumn{1}{l|}{198}            & \multicolumn{1}{l|}{187}            & 1107           \\ \hline
      \textbf{Female} & \multicolumn{1}{l|}{359}           & \multicolumn{1}{l|}{290}            & \multicolumn{1}{l|}{110}            & \multicolumn{1}{l|}{160}            & 919            \\ \hline
      \textbf{Total}  & \multicolumn{1}{l|}{729}           & \multicolumn{1}{l|}{642}            & \multicolumn{1}{l|}{308}            & \multicolumn{1}{l|}{347}            & 2026           \\ \hline
    \end{tabular}
  \end{table}


  Suppose that one of these 2026 students is randomly selected. Let $F$ denote the event that the
  selected individual is female, and $A,B,C$ and $D$ represent the events that they have blue, brown,
  green, and hazel eyes, respectively.
  \begin{enumerate}[(a)]
    \item Calculate $\P\left[F \right]$ and $\P\left[C \right]$
    \item Calculate $\P\left[F\cap C \right]$. Are these events independent? Why or why not?
    \item If the selected individual has green eyes, what's the probability that he or she is a female?
    \item If the selected individual is female, what is the probability that she has green eyes?
    \item What is the ``conditional distribution'' of eye color for females? (i.e., $\P\left[A|F
    \right], \P\left[B|F \right], \P\left[C|F \right], \P\left[D|F \right]$), and what is it for males?
    Compare the two distributions.
  \end{enumerate}
\end{problem}

\begin{problem}
  Two cards are randomly drawn from a deck of cards. Let
  \begin{align*}
    A&= [\text{at least one ace is drawn}]\\
    B&= [\text{both cards are aces}]\\
    S&= [\text{the ace of spades is drawn}]
  \end{align*}
  Compute the following conditional probabilities:
  \begin{enumerate}[(a)]
    \item $\P\left[S\mid B \right]$
    \item \label{item:3} $\P\left[B\mid S \right]$ % 1/17
    \item \label{item:4} $\P\left[B\mid A \right]$ % 1/33    
    \item Note that the answer to \ref{item:3} is almost twice as large as the
    answer to \ref{item:4}. Do you find this result surprising?
  \end{enumerate}
\end{problem}

\begin{problem}
  Recall the following definition:
  \begin{center}
    \fbox{Events $A$ and $B$ are \demph{\textbf{independent}} if and only if
      $\P\left[A\cap B \right] = \P\left[A \right]\P\left[B \right]$}
  \end{center}
  Assume that $A$ and $B$ are independent. Explain way $A^{c}$ and $B$ must be
  independent, and also why $A^{c}$ and $B^{c}$ must be independent.
\end{problem}


\newpage



\begin{problem}[Exercise 2.5.71]
  A space mining company currently has two active projects, one on the moon and one
  in the Kuiper belt. Let $L$ be the event that the lunar mining project is
  successful, and let $K$ be the event that the Kuiper belt mining project is successful. Suppose that $L$ and
  $K$ are independent events with $\P\left[L \right]= 0.4$ and $\P\left[K \right]= 0.7$.
  \begin{enumerate}[(a)]
    \item If the lunar project is not successful, what is the probability that
    the Kuiper belt project is also not successful? Explain your reasoning.
    \item What is the probability that at least one of the two mining projects will be successful?
    \item Given that at least one of the two projects is successful, what is the probability that only
    the lunar project is successful?
  \end{enumerate}
\end{problem}

\begin{problem}[Exercise 2.5.77]
  Five explosive devices are placed underneath the walls of the enemy
  stronghold, and are set to detonate simultaneously. Each bomb has a 4\%
  chance of failing to detonate at the specified time. Assuming that the
  probabilities are independent, calculate
  $\P\left[\text{at least one detonation occurs} \right]$ and
  $\P\left[\text{at least one bomb fails to detonate} \right]$
\end{problem}

\begin{problem}[Similar to Example 2.31]
  An email provider observes that $50\%$ of emails are spam. It intends to filter spam emails before they
  reach the inbox. The filter's accuracy for detecting spam email is $99\%$ and chances of tagging a
  non-spam email as spam is $5\%$.
  \begin{enumerate}[(a)]
    \item Draw a tree diagram to illustrate this situation
    \item If a certain email is tagged as spam, find the probability that it is not a spam email. That
    is, find
    \begin{equation*}
      \P\left[\text{email is not spam}\mid \text{email is tagged as spam} \right]
    \end{equation*}
  \end{enumerate}
\end{problem}



\begin{problem}
  Consider a universe where Netflix has exactly 6 shows, numbered $1$ through
  $6$. Suppose we randomly select a subscriber and ask which shows they watch.
  For each $i=1,2,\ldots,6$, let $E_{i}= \left\{i\right\} $ be the event that
  the subscriber watches show $i$. Suppose that
  \begin{equation*}
    P(E_1) = P(E_6) =  .10,  \ \ P(E_2) = P(E_5) = .15, \ \  P(E_3) = P(E_4) = .25.
  \end{equation*}
  Define events $A, B, C$ by
  \begin{equation*}
    A = \{2, 4, 6\},\ \  B = \{1, 2, 3\}, \ \ C = \{2, 3, 4, 5\}, \ \ D = \{1, 3, 5\}.
  \end{equation*}
  \begin{enumerate}[(a)]
    \item Are $A$ and $B$ dependent or independent?
    \item Are $A$ and $C$ dependent or independent?
    \item Are $A$ and $D$ dependent?
    \item Can two mutually exclusive events $A$ and $B$ with $P(A)>0$ and $P(B)>0$ be independent? Why or
    why not?
  \end{enumerate}
\end{problem}


\begin{problem}%[conditional probability]
  The prevalence of bird flu in a wild bird population is $1\%$. A new diagnostic test is advertised as
  having a false negative rate of $0.03$ and a false positive rate of $0.05$. To be precise, this means
  that
  \begin{equation*}
    \P\left[\text{test is negative}\mid \text{bird has the flu} \right] = 0.03
  \end{equation*}
  \begin{equation*}
    \P\left[\text{test is positive}\mid \text{bird doesn't have the flu} \right] = 0.05
  \end{equation*}
  \begin{enumerate}[(a)]
    \item What is the probability that a randomly-selected bird tests positive? \textit{(You don't need to
      simplify your answer.)}
    \item Suppose that a randomly-selected bird tests positive. What is the probability that the bird has
    the disease? \textit{(You don't need to simplify your answer.)}
  \end{enumerate}
\end{problem}






\newpage
\begin{problem}[Rolling two dice]
  Two 6-sided dice are rolled and added together. Let $X$ be the sum. 
  \begin{enumerate}[(a)]
    \item What is the probability mass function of $X$? [\textit{Hint:} First determine $p(1),p(2),\ldots$ and so on.]
    \item Determine the cumulative density function of $X$ and graph it.
    \item What is the expected value of $X$?
  \end{enumerate}
\end{problem}


\begin{problem}[Rolling a d6 `with advantage']
  In the game \textit{Dungeons \& Dragons}, there is a dice-rolling mechanic called ``rolling with
  advantage'' in which the player gets to roll a dice twice and keep the largest value. This bonus is
  typically regarded as a good thing, as higher rolls tend to be better in the game. In this problem,
  we'll investigate how significant this bonus is.
  % 
  Two 6-sided dice are rolled. Let $M$ be the maximum of the two rolls.
  \begin{enumerate}[(a)]
    \item What is the probability mass function of $M$? [\textit{Hint:} First determine $p(1),p(2),\ldots$ and so on.]
    \item Determine the cumulative density function of $M$ and graph it.
    \item What is the expected value of $M$?
    \item Suppose you are given the choice between two magical weapons: a magical sword that deals
    $1\text{d}6+1$ damage, and a magical spear that deals $1\text{d}6$ damage \textit{with advantage}.
    Which magical weapons has a higher expected damage? \textit{(Here, 1d6 means ``roll one six-sided
      dice''.)}
  \end{enumerate}
\end{problem}



\end{document}

